% !TeX root = ../papers/nc-p3.tex
\section{任意の正則辺における純関係の正次元成分}
\label{sec:pure-relation-curve-supports}
\begin{theorem}\label{thm:pure-relation-curve-supports}
標準正則な辺 $R\subset U\otimes V$ の純関係の被約軌跡を考える。
正次元なら次の二型に限る。
\begin{enumerate}
\item 一方の因子が固定された直線があり、そのような直線は各向き高々一本である。
正次元部分はこれらの直線の和である。
\item 二つの射影直線間の射影同型のグラフである滑らかな平面二次曲線が一本ある。
その外の純関係は高々三点であり、他の正次元成分はない。
\end{enumerate}
特に純関係の曲面は存在しない。これは元の辺の中の分解可能テンソルについての制限で、
辺の点対応（関係の零点）とは区別する。
\end{theorem}
\begin{proof}[部分空間に支持される関係数]
$D\subset U$ に対して $k_D=\dim(R\cap(D\otimes V))$ とする。
$\dim D=1,2,3$ に対して上界は順に $2,3,4$ である。
一次元の場合は定理\ref{thm:branching-centers-curve}の最初の評価、
三次元の場合は法線での端支持の下界2による。
二次元で $k_D\ge4$ なら $D^\perp$ の全点で評価階数は高々2、従って常に2である。
しかし定理\ref{thm:plane-branch-pencils}によれば、この分岐直線への制限に残る関係は
二次元以下でなく丁度三次元である。矛盾によって上界3を得る。右支持にも同じ上界がある。

純関係の整曲線の両支持の次元を $a,b$、そのテンソルの包の次元を $c$ とする。
関数体の積空間の付値から $c\ge a+b-1$ である。両方が4なら $c\ge7$ は不可能である。
一方の支持が真なら $c\le4$、二次元支持なら $c\le3$ である。
従って両方が動く場合は $a=b=2,c=3$ しかない。固定二平面のSegre二次曲面の平面切断で、
整かつ両因子が動くため滑らかな二次曲線である。片方が固定の場合は $k_D\le2$ により直線である。
二次曲線の三空間 $P$ は、両端の二平面 $A,B$ を同定すれば対称行列全体である。
\end{proof}
\begin{proof}[複数の曲線と孤立した純関係]
同じ向きの異なる二直線の第一因子が同じなら $k_D\le2$ に反し、異なるなら
その二因子を張る二平面に四次元の関係が支持される。従って同じ向きには一本しかない。
二次曲線と純直線の共存も不可能である。固定因子が $A$ 内なら三空間 $P$ の外の関係を
$A$ 内に生じ、外なら $A$ とその固定因子を張る三平面に五次元の関係を生じる。
右側でも同様である。

二つの二次曲線の左支持を $A,A'$ とする。同じなら両者とも同じ三空間 $P$ に入り、
そこにある二次曲線は一つである。交わりが一次元なら二曲線の包の交叉は高々一次元である。
実際、各包のうちその固定第一因子に支持されるのは対応する一点の直線のみである。
よって $A+A'$ の三平面に五次元以上の関係が入り矛盾する。
交わりが零なら $R$ 全体が両三空間の直和となる。$A^\perp$ のどの非零点で評価しても
$A'$ の対称二行列の評価だけが残り、その像は常に第二の曲線の同じ右二平面である。
これは二次元端支持を定める点の一意性に反する。以上で曲線の個数が決まり、
曲面なら無限個の異なる整曲線を含むので除かれる。

二次曲線がある場合に $U$ を $A$ で割った像 $R'$ は三次元である。
二次曲線の外の純点 $x\otimes y$ は $x\notin A,y\notin B$ を満たす。
例えば $y\in B$ なら $B$ に支持された全関係は $P$ に等しいためである。
$\bar x\otimes y\in R'$ はその純点を一意に決める。持上げの差は $A\otimes ky$ と
$P\subset A\otimes B$ の双方に入るからである。
$R'$ の要素を $U/A$ の双対の非零点で評価する四行三列の一次式鉛筆を取る。
全点で階数は2以上である。一般階数3なら三次小行列の共通零点は高々三点で、
階数2の各点にただ一つの核線があり、これが $R'$ の対応する純関係である。
一般階数2なら $A^\perp$ は分岐中心の直線になる。定理\ref{thm:plane-branch-pencils}の
三関係表示で、$R'$ の純関係は全て同じ第二因子 $b$ を持つ。
実際第三関係の二つの係数は $b$ を法として独立である。
従って二つの異なる純点が持ち上がれば $R$ に $b$ を持つ純直線を生じてしまう。
この場合に持ち上がるのは高々一点である。
\end{proof}
\path{computations/verify_pure_relation_curve_supports.py}は両支持の数値条件と三次鉛筆の次数を検算する。
